\documentclass[10pt,a4paper]{amsart}
\usepackage[margin=19mm]{geometry}
\usepackage{amsmath,amssymb,mathtools}
\usepackage[scr=rsfs,cal=euler]{mathalfa}
\usepackage{xfrac}
\usepackage[unicode,hidelinks,bookmarks=false]{hyperref}
\newcommand{\ie}{i.e.\ }
\newcommand{\cf}{cf.\ }
\newcommand{\resp}{respectively\ }
\newcommand{\Subsection}{\subsection}
\providecommand{\nobreakdash}{\nobreak\hbox{-}}
\setlength{\emergencystretch}{2em}
\multlinegap0pt
\usepackage{bm}
\usepackage{bbm}
\usepackage{dsfont}
\usepackage{xspace}
\usepackage{cancel}
\usepackage{mathtools}
\usepackage[shortlabels]{enumitem}
\usepackage{amsthm}
\makeatletter
\@ifundefined{definition}{\newtheorem{definition}{Definition}}{}
\@ifundefined{lemma}{\newtheorem{lemma}[definition]{Lemma}}{}
\@ifundefined{proposition}{\newtheorem{proposition}[definition]{Proposition}}{}
\@ifundefined{theorem}{\newtheorem{theorem}[definition]{Theorem}}{}
\makeatother
\def\R{{\mathbb{R}}}
\newcommand{\spt}{\mathrm{spt}\,}
\title[Generic mean curvature flow with obstacles]{Generic mean curvature flow with obstacles}
\author{Tim Laux, Keisuke Takasao}
\date{Source: arXiv:2507.06150v1 (CC BY 4.0)}
\begin{document}
\maketitle

\noindent\emph{Source and licence.} Generic mean curvature flow with obstacles. Version arXiv:2507.06150v1 (CC BY 4.0), distributed under the Creative Commons Attribution 4.0 International licence, as recorded in the arXiv licence metadata for this exact version and shown on the abstract page https://arxiv.org/abs/2507.06150v1. Full rights notice: licenses/arxiv-2507.06150-CC-BY-4.0.txt.\par\smallskip

\noindent\textit{Editorial scope.} Counted unit: the Theorem whose source label is \texttt{thm:3.4} in the article (source0.tex lines 1595--1605, source label \texttt{thm:3.4}), delivered together with the article's own complete original proof of it (source0.tex lines 1607--1776, one \texttt{proof} environment of 170 lines). No mathematics is written, repaired or re-derived: statement and proof are the article's own text line for line. Required context retained uncounted: def:well-prepared (lines 285--313); lem3.4:assumption1 (lines 293--301); prop:mp (lines 635--669); lemma3.3 (lines 1415--1439); BVineq (lines 1431--1437). Every definition, display and label of those lines is retained unchanged. Omitted from this package and not redistributed: 1--284, 302--634, 670--1414, 1438--1594, 1606--1606, 1777--2727. Nothing inside a retained excerpt is omitted or altered: every retained body is delivered exactly as the article's lines are written, so no omission receipt of any kind accompanies this package. Citations: the retained excerpts carry no citation command at all, so no bibliography entry of the article is transcribed and no reference is redistributed. Reference adaptation: none was needed and none was made; every reference inside the delivered unit resolves inside it, so no unresolved reference or citation is produced. Printed slips kept literal: no printed word, symbol, hypothesis or number of the counted statement or of its proof was altered. Layout and class: the article's own class is \texttt{amsart} with packages including todonotes, enumerate, babel, geometry, amssymb, latexsym, amsfonts, mathrsfs, amsmath, amsthm, multicol, constants, color, bm; the wrapper is the lane's pinned \texttt{amsart} editorial wrapper at 10pt on A4, which restates the theorem-like environments the retained text needs and adds the article's own macro definitions that the retained text needs (\texttt{\textbackslash R}, \texttt{\textbackslash spt}), with extra wrapper packages (bm, bbm, dsfont, xspace, cancel, mathtools, enumitem). The delivered numbering is the unit's own. Assets: no retained span contains a figure environment, a graphics-inclusion command, a PDF-page inclusion, a table or a TikZ picture; no image file is copied into this package, the package declares no asset dependency and no image file is redistributed with it. Licence: version arXiv:2507.06150v1 of this article is distributed under the Creative Commons Attribution 4.0 International licence, as recorded in the arXiv licence metadata for this exact version and shown on the abstract page https://arxiv.org/abs/2507.06150v1. A full rights notice is delivered at licenses/arxiv-2507.06150-CC-BY-4.0.txt. Characters: every retained excerpt is pure ASCII, so the delivered engine, XeLaTeX with its default Latin Modern text font, reports no missing character.
\begin{definition}[Well-prepared data]\label{def:well-prepared}
    We call $\phi,\psi,g\colon \Omega \to \R$ \emph{well-prepared data} if $\phi,\psi \in C^1 (\Omega), g \in C^2 (\Omega)$,
    \[
        \phi<\psi 
        \quad \text{and} \quad
        \phi \leq g \leq \psi \quad \text{in }\Omega,
    \]
    and there exists $L, \ell<\infty$ such that
    \begin{equation}\label{lem3.4:assumption1}
        \max_{x \in \Omega} |g(x)| \leq L,
        \quad 
        \max_{x \in \Omega} \phi (x) \leq L,
        \quad 
        \text{and} 
        \quad
        \min_{x \in \Omega} \psi (x) \geq -L.
    \end{equation}
    and the sets of critical points
    \[
        \{ x \in \Omega \mid \phi (x) \geq - (L+l) \quad \text{and} \quad
    |\nabla \phi (x)|=0 \}
    \]
    and
    \[
    \{ x \in \Omega \mid \psi (x) \leq L+l \quad \text{and} \quad
    |\nabla \psi (x)|=0 \}
    \]
    are finite.
\end{definition}

\begin{proposition}\label{prop:mp}
For any $T>0$ we have the following a priori estimates.
\begin{enumerate}[(i)]
    \item \label{item:mp1}
    Uniform bounds:
    \[
    \min \{  
    \min_{x \in \Omega} \psi (x), \min_{x \in \Omega } g (x) \}
    \leq
    \min _{(x,t) \in \Omega_T } u_\varepsilon (x,t) \leq 
    \max _{(x,t) \in \Omega_T } u_\varepsilon (x,t)
    \leq \max \{  
    \max_{x \in \Omega} \phi (x), \max_{x \in \Omega } g (x) \}.
    \]
    \item\label{item:mp2}
    Gradient bound:
    \[
    \max _{(x,t) \in \Omega_T } 
    | \nabla u_\varepsilon (x,t)|
    \leq \max \{  
    \max_{x \in \Omega} |\nabla \phi (x)|,
    \max_{x \in \Omega} |\nabla \psi (x)|,    
    \max_{x \in \Omega } |\nabla g (x)| \}.
    \]
    \item\label{item:mp3}
    Bound on time derivative:
    \[
    \max _{(x,t) \in \Omega_T } 
    | \partial_t u_\varepsilon (x,t)|
    \leq
    \max_{x \in \Omega } |\nabla ^2 g (x)|.
    \]
\end{enumerate}

\end{proposition}

\begin{lemma}\label{lemma3.3}
\begin{enumerate}[(i)]
\item \label{item:lemma3.3(i)} For any $\delta >0$ there exist $\varepsilon_0 >0$ and $\alpha>0$ such that
\begin{equation}\label{lemma3.3(i)}
u_\varepsilon (x,t) - \phi (x) \geq \alpha \quad \text{for any} \ (x,t) \in B_ - ^\delta
\quad \text{and} \quad
\psi (x) - u_\varepsilon (x,t) \geq \alpha \quad \text{for any} \ (x,t) \in B_ + ^\delta
\end{equation}
for any $\varepsilon \in (0,\varepsilon_0)$.
\item\label{item:lemma3.3(ii)} Let $\delta >0$ and $X \in C^1 (\Omega \times [0,T];\R^d)$ satisfy
\begin{equation}
    X \cdot \nabla \phi \leq 0 \quad \text{on} \ A_ - ^\delta
    \quad \text{and} \quad
    X \cdot \nabla \psi \geq 0 \quad \text{on} \ A_ + ^\delta
\end{equation}
Then for a.e.\ $t \in [0,T]$ we have
\begin{equation}\label{BVineq}
\int_{\Omega} \left\{
 \left( \delta_{ij} - \frac{ \partial_{x_i} u \partial _{x_j} u }{|\nabla u|^2} \right)
\partial_{x_j} X^i
+V\Big(\frac{-\nabla u}{|\nabla u|}\Big)  \cdot X\right\} |\nabla u| \, dx
\geq 0.
\end{equation}
\end{enumerate}
\end{lemma}

\begin{theorem}\label{thm:3.4}
Suppose $\phi, \psi, g$ are well-prepared according to Definition~\ref{def:well-prepared} and let $X \in C^1 (\Omega \times [0,T])$ satisfy
    \[
    X \cdot \nabla \phi \leq 0 \quad \text{on} \ A_-
    \quad
    \text{and}
    \quad
    X\cdot \nabla \psi \geq 0 \quad \text{on} \ A_+.
    \]
Then \eqref{BVineq} holds.
\end{theorem}

\begin{proof}
First we assume that 
\begin{equation}\label{lem3.4:assumption2}
\spt X \cap \left( \cup_{i=1} ^N B_r (x_i) \cup \cup_{i=1} ^M B_r (y_i) \right)
\times [0,T] = \emptyset
\end{equation}
for some $r>0$, 
where 
\[
\{x_1,\dots,x_N\}
=\{ x \in \Omega \mid \psi (x) \leq L+l \quad \text{and} \quad
    |\nabla \psi (x)|=0 \}
\]
and
\[
\{y_1,\dots,y_M \}
=\{ x \in \Omega \mid \phi (x) \geq -(L+l) \quad \text{and} \quad
    |\nabla \phi (x)|=0 \}.
\]
By the assumptions for $\psi$ and $\phi$, one can check that 
for $r>0$, there exists $C_r >0$ such that
\begin{equation}\label{lem3.4-0}
|\nabla \psi | \geq C_r \quad \text{on} \ \{ \psi (x) \leq L+l \} \setminus \cup_{i=1} ^N B_r (x_i)
\quad \text{and} \quad 
|\nabla \phi | \geq C_r \quad \text{on} \ \{ \phi (x) \geq -(L+l) \} \setminus \cup_{i=1} ^M B_r (y_i).
\end{equation}
By the maximum principle (Proposition \ref{prop:mp}) and \eqref{lem3.4:assumption1},
it holds that $\sup_{(x,t) \in \Omega [0,\infty)} |u_\varepsilon (x,t)| \leq L $ for any $\varepsilon >0$. Therefore 
\begin{equation*}
A_+ \subset \{ \psi \leq L \} \times [0,T]
\qquad \text{and} \qquad
A_- \subset \{ \phi \geq -L \} \times [0,T] .
\end{equation*}
Hence for any $\delta \in (0,l)$, 
\begin{equation}\label{lem3.4-1}
A_+ ^\delta \subset \{ \psi \leq L +l \} \times [0,T]
\qquad \text{and} \qquad
A_- ^\delta \subset \{ \phi \geq -(L+l) \} \times [0,T] .
\end{equation}
By \eqref{lem3.4-0} and \eqref{lem3.4-1}, if $(x,t) \in A_+ ^\delta$ and $|\nabla \psi (x)| <C_r$,
then $x \in \cup_{i=1} ^N B_r (x_i)$. Therefore \eqref{lem3.4:assumption2} implies that
\begin{equation}\label{lem3.4-2}
(x,t) \in \spt X \cap A_+^\delta
\quad \Rightarrow \quad
|\nabla \psi (x) | \geq C_r
\quad
(\text{resp.} \ (x,t) \in \spt X \cap A_- ^\delta
\quad \Rightarrow \quad
|\nabla \phi (x) | \geq C_r
)
\end{equation}
Set $\omega _+ (\delta) := \min_{A_+ ^\delta} (X \cdot \nabla \psi) \leq 0$ and
$\omega _- (\delta) := \max_{A_- ^\delta} (X \cdot \nabla \phi) \geq 0$. Note that
\begin{equation}
\omega _{\pm} (\delta) \to 0 \qquad \text{as} \ \delta \to 0,
\end{equation}
by the assumption of $X$. 
Due to \eqref{lem3.4-2}, we can choose $Y_\psi, Y_\phi \in C^1 (\Omega)$
depending only on $C_r$, $\psi$, and $\phi$ such that
\begin{equation}
\nabla \psi \cdot Y_\psi \geq \frac{C_r ^2}{2} \quad \text{on} \ \spt X \cap A_+ ^\delta 
\quad \text{and} \quad
\nabla \phi \cdot Y_\phi \geq \frac{C_r ^2}{2} \quad \text{on} \ \spt X \cap A_- ^\delta
\end{equation}
for any $\delta \in (0,l)$. 
Note that we may assume that
\[
\nabla \psi \cdot Y_\psi \geq 0 \quad \text{on} \ A_+ ^\delta 
\quad \text{and} \quad
\nabla \phi \cdot Y_\phi \geq 0 \quad \text{on} \ A_- ^\delta
\]
by using some suitable cut-off function for $\spt X \cap A_\pm ^\delta$.
Choose $\delta_\ast \in (0,l)$ such that $A_+ ^{2\delta_\ast} \cap A_- ^{2\delta_\ast} =\emptyset$ and let $\eta \in C^1 (\Omega)$ be a cut-off function for $A_+ ^{\delta_\ast}$, namely,
$0\leq \eta (x)\leq 1$ for any $x\in \Omega$,
$\eta =1$ on $A_+ ^{\delta_\ast}$, and $\eta =0$ on $A_- ^{\delta _\ast}$.
Set
\[
X_\delta := X - \frac{2}{C_r ^2} (\omega_+ (\delta) \eta Y_\psi 
+ \omega _- (\delta) (1-\eta) Y_\phi) \in C^1 (\Omega \times [0,T]) \quad \text{for any} \ \delta \in (0,\delta_\ast).
\]
If $(x,t) \in A_+ ^\delta \cap \spt X$, then 
\[
X_\delta \cdot \nabla \psi 
=
\left( X - \frac{2}{C_r ^2} \omega_+ (\delta) Y_\psi \right) \cdot \nabla \psi
\geq 
\omega_+ (\delta) - \frac{2}{C_r ^2} \omega_+ (\delta) \cdot \frac{C_r ^2}{2}\geq 0.
\]
In addition, $X_\delta \cdot \nabla \psi \geq 0$ on $A_+ ^\delta \cap (\spt X)^c$. Hence
$X_\delta \cdot \nabla \psi \geq 0$ on $A_+ ^\delta$ 
(resp. $X_\delta \cdot \nabla \phi \leq 0$ on $A_- ^\delta$)
and thus $X_\delta$ satisfies 
\begin{equation*}
\int_{\Omega} \left\{ 
\left( \delta_{ij} - \frac{ \partial_{x_i} u \partial _{x_j} u }{|\nabla u|^2} \right)
\partial_{x_j} X_\delta ^i 
+
X_\delta \cdot \left(\frac{-\nabla u}{|\nabla u|}\right) V 
\right\}
|\nabla u| \, dx \geq 0
\end{equation*}
for a.e.\ $t \in [0,T]$, by Lemma \ref{lemma3.3}. Since $\omega _\pm (\delta)\to 0$ as 
$\delta \to 0$, $X_\delta \to X$ in $C^1 (\Omega \times [0,T])$. 
Therefore \eqref{BVineq} holds 
for a.e.\ $t \in [0,T]$.

Next we prove \eqref{BVineq} without \eqref{lem3.4:assumption2}.
Let $r>0$ and $\eta_r \in C^1 (\Omega)$ be a cut-off function such that
$\eta_r =1$ on $\Omega \setminus (\cup_{i=1} ^N B_{2r} (x_i) \cup \cup_{i=1} ^M B_{2r} (y_i))$
and $\eta_r=0$ on $\cup_{i=1} ^N B_{r} (x_i) \cup \cup_{i=1} ^M B_{r} (y_i)$.
Then 
\begin{equation}\label{eta-r1}
\int_{\Omega} \left\{
\left( \delta_{ij} - \frac{ \partial_{x_i} u \partial _{x_j} u }{|\nabla u|^2} \right)
\partial_{x_j} (\eta_r X ^i) 
+\eta_r X \cdot \left( \frac{-\nabla u}{|\nabla u|} \right) V
\right\} |\nabla u| \, dx \geq 0
\end{equation}
for a.e.\ $t \in [0,T]$, by the first step. Since $\eta_r \to 1$ a.e.\ $\Omega$ as $r\to 0$, 
for a.e.\ $t \in [0,T]$ we have
\begin{equation}\label{eta-r2}
\int_{\Omega} \eta_r X \cdot \left( \frac{-\nabla u}{|\nabla u|} \right) V |\nabla u| \, dx
\to
\int_{\Omega} X \cdot \left(\frac{-\nabla u}{|\nabla u|} \right) V |\nabla u| \, dx \qquad \text{as} \ r\to 0
\end{equation}
by the dominated convergence theorem. Note that $| V| |\nabla u| =|\partial _t u | $ a.e.\ $(x,t)$ and $\| \partial _t u (\cdot,t) \|_{L^\infty (\Omega)} <\infty$ for a.e.\ $t \geq 0$. 
We compute that
\begin{equation*}
\begin{split}
\int_{\Omega} |\nabla \eta_r | \, dx
= & \, \int_0 ^1 \mathcal{H}^{d-1} (\{ \eta_r =s \}) \, ds
\leq \int_0 ^1 \mathcal{H}^{d-1} (\cup_{i=1} ^N \partial B_{2r} (x_i) \cup \cup_{i=1} ^M \partial B_{2r} (y_i)) \, ds \\
\leq & \, \omega _{d-1} (M+N) (2r)^{d-1},
\end{split}
\end{equation*}
where we used the coarea formula and $\omega_{d-1} := \mathcal{H}^{d-1} (\partial B_1 (0))$.
Therefore
\begin{equation*}
\begin{split}
&\left| \int_{\Omega} \left( \delta_{ij} - \frac{ \partial_{x_i} u \partial _{x_j} u }{|\nabla u|^2} \right)
(\partial_{x_j} \eta_r) X ^i |\nabla u| \, dx \right|
=
\left| \int_{\Omega} \left( \frac{\nabla u}{|\nabla u|} \cdot \nabla \eta_r \right) 
\left( \frac{\nabla u}{|\nabla u|} \cdot X \right)  |\nabla u| \, dx \right| \\
\leq & \, 
\max _{\Omega\times [0,T]} |X| \| \nabla u \|_{L^\infty (\Omega \times [0,T])}
\int_{\Omega} |\nabla \eta_r | \, dx
\to 0 \qquad \text{as} \ r\to 0
\end{split}
\end{equation*}
and thus
\begin{equation}\label{eta-r3}
\begin{split}
&\int_{\Omega} \left( \delta_{ij} - \frac{ \partial_{x_i} u \partial _{x_j} u }{|\nabla u|^2} \right)
\partial_{x_j} (\eta_r X ^i) |\nabla u| \, dx \\
= & \, 
\int_{\Omega} \left( \delta_{ij} - \frac{ \partial_{x_i} u \partial _{x_j} u }{|\nabla u|^2} \right)
(\partial_{x_j} \eta_r) X ^i |\nabla u| \, dx
+
\int_{\Omega} \eta_r 
\left( \delta_{ij} - \frac{ \partial_{x_i} u \partial _{x_j} u }{|\nabla u|^2} \right)
\partial_{x_j} X ^i |\nabla u| \, dx \\
\to & 
\int_{\Omega} \left( \delta_{ij} - \frac{ \partial_{x_i} u \partial _{x_j} u }{|\nabla u|^2} \right)
\partial_{x_j} X ^i |\nabla u| \, dx \qquad \text{as} \ r\to 0,
\end{split}
\end{equation}
where we used the dominated convergence theorem again.
By \eqref{eta-r1}, \eqref{eta-r2}, and \eqref{eta-r3} we obtain \eqref{BVineq}.
\end{proof}

\end{document}
